\subsection{Replay Benchmarks and Out-of-Order Stream Evaluation}

Replaying recorded data through a streaming pipeline is a common way to evaluate stream processors. Existing work differs from SmokeWatch in what is replayed, what is controlled and what is checked.

\textbf{Engine benchmarks.} The Yahoo streaming benchmark~\cite{chintapalli2016benchmarking} compares Storm, Flink and Spark Streaming on a full Kafka and Redis pipeline with a generated advertising workload, and reports latency and throughput. Karimov et al.~\cite{karimov2018benchmarking} separate the system under test from the data driver and define sustainable throughput for windowed operators, using workloads inspired by online gaming. ESPBench~\cite{hesse2021} adds a toolkit for data ingestion, query result validation and automated runs, and evaluates Spark, Flink and Hazelcast Jet. Result validation is therefore not new to our project. These benchmarks compare engines on enterprise or generated workloads, whereas SmokeWatch fixes the engine (Kafka and Spark) and compares a replayed environmental workload with a batch reference. Our workload is small (18{,}000 events), so we make no claim about sustainable throughput or engine ranking.

\textbf{Deterministic delay.} Grulich et al.~\cite{grulich2019generating} already provide a deterministic, seeded generator that replays prerecorded files with configurable out-of-order fractions and delays, so deterministic delay is also not new, and we do not claim a new delay generator. SmokeWatch's delay profiles are a simpler implementation choice for our workload: the delay of an event is derived from its event identifier and a named profile (none, moderate or severe), so the same events are delayed in every run, and the arrival time is the event hour plus that delay. The scenario label does not change which events are delayed. Each Kafka message also carries a \texttt{run\_code}, so storage can deduplicate by run and event and keep runs separate. The workload differs as well: three sources (NASA FIRMS, EPA PM2.5 and NOAA weather) are replayed together for one fire event.

\textbf{Accelerated replay.} ReStream~\cite{schleiersmith2016} replays stored event logs much faster than real time for backtesting, and focuses on obtaining the same results as serial processing in log order. SmokeWatch does not build a new engine. It uses standard Kafka and Spark at fixed replay speeds, so our contribution is the experimental design rather than a faster replay mechanism.

\textbf{Event time and late data.} The Dataflow model~\cite{akidau2015} and later work on watermarks~\cite{akidau2021} give the vocabulary we rely on: event time versus processing time, and watermarks as a statement of input completeness. Our delay profiles create a gap between event time and arrival time on real data. The freshness values reported for our runs are computed from this scheduled arrival delay, not from observed end-to-end completion time, so they describe the injected delay rather than measured pipeline latency.

\textbf{Contribution and status.} SmokeWatch does not propose a new replay algorithm or benchmark, and we make no claim of benchmark superiority. Its contribution is a small replay design for real multi-source environmental data that combines deterministic delay profiles, run codes carried from Kafka messages into storage, and a batch reference for comparison. Implemented so far are the replay producer with three delay profiles, run codes in every message, and the Kafka, Spark and PostgreSQL path exercised on diagnostic runs. Still pending are the official-data integrated runs on the checksum-verified inputs, the row-level comparison with the batch reference, and MongoDB ingestion in the same run, which is not yet confirmed. Results are limited to one fixture size and one fire event.
